
\Slide{1}{课程与学习目标}
这门课用数学语言研究经济均衡。“存在”问的是满足全部条件的对象能否找到，不一定要求算出它。先掌握集合、逻辑和极限，再学习连续性、凸性和不动点；不必提前学习微积分。证明题的基本格式是：写假设，写目标，选任意对象，逐步调用定义，最后回到目标。

\Slide{2}{参考文献如何使用}
四项文献是课程的知识来源，不是四份额外作业。本讲义以课件为主，用 Debreu (1959) 第一章核对术语。Berge 对应连续性和最优化工具，后续章节使用这些工具研究经济均衡。不要把阅读历史书目等同于证明训练。

\Slide{3}{教材七章的逻辑}
商品和价格把经济对象表示为向量；生产者和消费者描述可行集合与偏好；均衡要求各方选择相容；最优讨论效率；不确定性扩大商品的描述。这十二讲主要负责第一章数学工具，后六章在此只作动机。

\Slide{4}{两种均衡研究路径}
同时最优化强调各人的最优选择能同时实现；超额需求方法用“需求减供给”表示市场失衡。连续性和不动点是存在论证的桥梁。测度空间代理人属于更高级材料，本讲无需提前掌握测度论。

\Slide{5}{均衡思想的来源}
“看不见的手”是直觉性表述，数学均衡需要列出价格、需求、供给及市场出清等条件。Walras 和 Arrow--Debreu 的工作把问题写成可检验的数学命题。不能从历史口号直接推出任何经济都必有均衡，必须检查模型假设。

\Slide{6}{Adam Smith}
本页介绍人物、日期和著作。对证明考试，重点是理解经济问题的动机，不把人物年表当成数学推理。形式化之后，“某种价格能够协调选择”才成为待证明的存在命题。

\Slide{7}{Cournot}
Cournot 把相互影响的选择放进数学模型。一个简单动机是两家企业分别选择产量，各自利润取决于两家的选择。一个人的最优选择可能是集合，这解释了为什么后面除了函数还需要“对应”。

\Slide{8}{Walras}
一般均衡同时考虑多个市场。若价格是 $p$，把总需求与总供给之差记为 $z(p)$，市场出清可抽象成求 $z(p)=0$。这只是说明数学问题的形式；本页没有给出 $z$ 的具体模型或存在证明。

\Slide{9}{Edgeworth}
本页列出人物、著作和极限定理的名称。“极限”描述一族对象在某种变化下的趋近行为，后面从数列开始学习。不要凭名称猜测定理的假设与结论；考试引用结果时要写完整适用条件。

\Slide{10}{Pareto 与效率}
若一种可行改变使至少一人更好、无人更差，称为 Pareto 改进；不存在这种改进的可行配置称为 Pareto 最优。它不意味着人人同样富，也不自动意味着存在支持它的均衡价格，后者需要凸性和分离工具。

\Slide{11}{von Neumann}
极小极大思想研究两个决策者目标相反时的策略选择。最优化结果经常需要紧性保证最优点存在，凸性允许混合。这里先了解用途，等学完107、109再理解其数学基础。

\Slide{12}{Kakutani}
对应把一个输入送到一组可能输出。不动点的形式是 $x\in\Phi(x)$，意思是当前选择本身属于允许的响应集。110将给出非空、紧、凸和半连续等条件，不能只背“不动点一定存在”。

\Slide{13}{Nash}
均衡可以解释为每个人的选择都是对其他人选择的最佳响应。把所有最佳响应组合成对应，就能把均衡存在转为不动点问题。Kuhn--Tucker 条件属于优化的另一套工具；本页提到名称不等于要求先掌握它的证明。

\Slide{14}{Arrow}
本页涉及社会选择、一般均衡、福利、增长和信息等研究领域。不同领域的定理假设差别很大。学习方法是区分“某人的研究主题”和“本课程已建立的数学命题”，不从人物介绍推出未经证明的结论。

\Slide{15}{Debreu}
本课采用其公理化思路：先明确对象与性质，再从假设推出结论。例：要证明连续函数在集合上取最大值，必须补充集合非空且紧等条件；只写“函数连续”不够。

\Slide{16}{人物年表汇总}
本页重复汇总十位人物，前页的历史信息都保留在原图中。建议把学习笔记的时间用于记录“这个概念用来解决什么问题”，数学主线是：可行集合、最优选择、连续响应、均衡不动点。

\Slide{17}{另一份经济学家名单}
本页为历史名单。原文 David Richard 通常应为 David Ricardo（大卫·李嘉图）；这是人物名称说明，不影响后续数学。此页没有数学命题需要证明。

\Slide{18}{二十世纪人物}
本页突出 Samuelson 与 Arrow；“Great”属于课程介绍的评价性表述。我们的学习任务仍是理解严谨推导，而不是据此给研究者排序。没有新增数学假设、公式或证明。

\Slide{19}{四种存在性方法}
不动点方法是本组讲义的主要目标；算法方法研究近似解；指数和度理论研究映射的整体结构；微分过程研究演化。存在、唯一、可计算、收敛是不同问题。例如 $f(x)=x$ 在区间上有无穷多个不动点：存在并不推出唯一。
